% ============================================================
% macros.tex — Style System & Custom Commands
% ============================================================

% ======================================================================
% A. STYLE SYSTEM
% ======================================================================

% --- Document info setters ---
\let\DocTitle\empty
\let\DocSubtitle\empty
\let\DocContact\empty
\let\DocAuthor\empty
\let\DocDate\empty
\let\DocHeader\empty

\newcommand{\doctitle}[1]{\def\DocTitle{#1}}
\newcommand{\docsubtitle}[1]{\def\DocSubtitle{#1}}
\newcommand{\doccontact}[1]{\def\DocContact{#1}}
\newcommand{\docauthor}[1]{\def\DocAuthor{#1}}
\newcommand{\docdate}[1]{\def\DocDate{#1}}
\newcommand{\docheader}[1]{\def\DocHeader{#1}}

\makeatletter

% --- Style dispatcher ---
\newcommand{\usestyle}[1]{%
    \ifcsname style@#1\endcsname
        \csname style@#1\endcsname
    \else
        \PackageError{macros}{Unknown style '#1'}{%
            Available styles: academic, homework, note, minimal, custom}%
    \fi
}

% ------------------------------------------------------------------
% Style: academic
%   Title: centered, large title + subtitle + author + date
%   Sections: \Large bold with underline, numbered
%   Header: empty | Footer: centered page number
% ------------------------------------------------------------------
\newcommand{\style@academic}{%
    % Title
    \renewcommand{\maketitle}{%
        \begin{center}
            {\fontsize{30pt}{36pt}\selectfont\bfseries\color{\TitleColor}\DocTitle\par}
            \vspace{6pt}
            \ifx\DocSubtitle\empty\else
                {\fontsize{12pt}{16pt}\selectfont\color{\HeadingColor}\DocSubtitle\par}
                \vspace{4pt}
            \fi
            \ifx\DocAuthor\empty\else
                {\normalsize\DocAuthor\par}
                \vspace{2pt}
            \fi
            \ifx\DocDate\empty\else
                {\small\color{gray}\DocDate\par}
                \vspace{2pt}
            \fi
            \ifx\DocContact\empty\else
                {\small\color{gray}\DocContact\par}
            \fi
        \end{center}
        \vspace{12pt}
    }%
    % Sections
    \titleformat{\section}
        {\color{\HeadingColor}\Large\bfseries}
        {\thesection}{1em}{}
        [\vspace{-4pt}\color{\HeadingColor}\titlerule]
    \titlespacing{\section}{0pt}{12pt}{2pt}
    \titleformat{\subsection}
        {\color{\HeadingColor}\large\bfseries}
        {\thesubsection}{1em}{}
    \titlespacing{\subsection}{0pt}{10pt}{0pt}
    \titleformat{\subsubsection}
        {\color{\HeadingColor}\normalsize\bfseries}
        {\thesubsubsection}{1em}{}
    \titlespacing{\subsubsection}{0pt}{6pt}{0pt}
    % Header / Footer
    \fancyhf{}
    \fancyfoot[C]{\thepage}
    \renewcommand{\headrulewidth}{0pt}
    \renewcommand{\footrulewidth}{0pt}
}

% ------------------------------------------------------------------
% Style: homework
%   Title: left-right aligned table (title+subtitle left, date+author right), hrule
%   Sections: \large bold, no underline, NO numbering (type anything in heading)
%   Header: left=\docheader (defaults to subtitle), right=author | Footer: centered page number
%   Header rule: 0.4pt
% ------------------------------------------------------------------
\newcommand{\style@homework}{%
    % Title (first page has no header)
    \renewcommand{\maketitle}{%
        \thispagestyle{plain}%
        \noindent
        \begin{tabular*}{\textwidth}{@{}l@{\extracolsep{\fill}}r@{}}
            {\fontsize{30pt}{36pt}\selectfont\bfseries\color{\TitleColor}\DocTitle}
            & {\normalsize\DocDate} \\[4pt]
            {\fontsize{12pt}{16pt}\selectfont\color{\HeadingColor}\DocSubtitle}
            & {\normalsize\DocAuthor} \\
        \end{tabular*}
        \ifx\DocContact\empty\else
            \\[2pt]\noindent{\small\color{gray}\DocContact}
        \fi
        \vspace{4pt}
        \hrule
        \vspace{12pt}
    }%
    % Sections: all unnumbered (type anything in heading); h1 has rule
    \titleformat{\section}
        {\color{\HeadingColor}\LARGE\bfseries}
        {}{0em}{}
        [\vspace{-4pt}\color{gray!70}\titlerule]
    \titlespacing{\section}{0pt}{10pt}{2pt}
    \titleformat{\subsection}
        {\color{\HeadingColor}\Large\bfseries}
        {}{0em}{}
    \titlespacing{\subsection}{0pt}{8pt}{-2pt}
    \titleformat{\subsubsection}
        {\color{\HeadingColor}\large\bfseries}
        {}{0em}{}
    \titlespacing{\subsubsection}{0pt}{6pt}{-2pt}
    % Header / Footer
    \fancyhf{}
    \fancyhead[L]{\small\ifx\DocHeader\empty\DocSubtitle\else\DocHeader\fi}
    \fancyhead[R]{\small\DocAuthor}
    \fancyfoot[C]{\small\thepage}
    \renewcommand{\headrulewidth}{0.4pt}
    \renewcommand{\footrulewidth}{0pt}
    \fancypagestyle{plain}{\fancyhf{}\fancyfoot[C]{\small\thepage}\renewcommand{\headrulewidth}{0pt}}%
}

% ------------------------------------------------------------------
% Style: note
%   Title: left-aligned 24pt title, subtitle, author · date inline, hrule
%   Sections: \LARGE bold with underline, numbered
%   Header: left=\docheader (defaults to title), right=author | Footer: centered page number
%   Header rule: 0.4pt
% ------------------------------------------------------------------
\newcommand{\style@note}{%
    % Title (first page has no header)
    \renewcommand{\maketitle}{%
        \thispagestyle{plain}%
        \noindent
        {\fontsize{24pt}{30pt}\selectfont\bfseries\color{\TitleColor}\DocTitle\par}
        \vspace{4pt}
        \ifx\DocSubtitle\empty\else
            {\vspace{-6pt}\fontsize{12pt}{16pt}\selectfont\color{\HeadingColor}\DocSubtitle\par}
        \fi
        \ifx\DocAuthor\empty\else
            {\vspace{-6pt}\normalsize\DocAuthor}%
            \ifx\DocDate\empty\else{ · \DocDate}\fi
            \par
        \fi
        \ifx\DocContact\empty\else
            {\small\color{gray}\DocContact\par}
        \fi
        \vspace{4pt}
        \hrule
        \vspace{12pt}
    }%
    % Sections: all unnumbered (type anything in heading); h1 has rule
    \titleformat{\section}
        {\color{\HeadingColor}\LARGE\bfseries}
        {\thesection.}{0.25em}{}
        [\vspace{-4pt}\color{gray!70}\titlerule]
    \titlespacing{\section}{0pt}{10pt}{2pt}
    \titleformat{\subsection}
        {\color{\HeadingColor}\Large\bfseries}
        {}{0em}{}
    \titlespacing{\subsection}{0pt}{8pt}{-2pt}
    \titleformat{\subsubsection}
        {\color{\HeadingColor}\large\bfseries}
        {}{0em}{}
    \titlespacing{\subsubsection}{0pt}{6pt}{-2pt}
    % Header / Footer
    \fancyhf{}
    \fancyhead[L]{\small\ifx\DocHeader\empty\DocTitle\else\DocHeader\fi}
    \fancyhead[R]{\small\DocAuthor}
    \fancyfoot[C]{\small\thepage}
    \renewcommand{\headrulewidth}{0.4pt}
    \renewcommand{\footrulewidth}{0pt}
    \fancypagestyle{plain}{\fancyhf{}\fancyfoot[C]{\small\thepage}\renewcommand{\headrulewidth}{0pt}}%
}

% ------------------------------------------------------------------
% Style: minimal
%   Title: left-aligned, subtitle · author · date inline
%   Sections: \large bold, no underline, NO numbering
%   Header: empty | Footer: centered page number
% ------------------------------------------------------------------
\newcommand{\style@minimal}{%
    % Title
    \renewcommand{\maketitle}{%
        \noindent
        {\fontsize{30pt}{36pt}\selectfont\bfseries\color{\TitleColor}\DocTitle\par}
        \vspace{4pt}
        \ifx\DocSubtitle\empty\else
            {\fontsize{12pt}{16pt}\selectfont\color{\HeadingColor}\DocSubtitle\par}
            \vspace{2pt}
        \fi
        \ifx\DocAuthor\empty\else
            {\normalsize\DocAuthor}%
            \ifx\DocDate\empty\else{ · \DocDate}\fi
            \par
        \fi
        \vspace{12pt}
    }%
    % Sections (no numbering)
    \titleformat{\section}
        {\color{\HeadingColor}\Large\bfseries}
        {}{0em}{}
    \titlespacing{\section}{0pt}{12pt}{2pt}
    \titleformat{\subsection}
        {\color{\HeadingColor}\large\bfseries}
        {}{0em}{}
    \titlespacing{\subsection}{0pt}{8pt}{0pt}
    \titleformat{\subsubsection}
        {\color{\HeadingColor}\normalsize\bfseries}
        {}{0em}{}
    \titlespacing{\subsubsection}{0pt}{6pt}{0pt}
    % Header / Footer
    \fancyhf{}
    \fancyfoot[C]{\thepage}
    \renewcommand{\headrulewidth}{0pt}
    \renewcommand{\footrulewidth}{0pt}
}

% ------------------------------------------------------------------
% Style: custom (academic as base; override with \renewcommand / \titleformat after)
% ------------------------------------------------------------------
\newcommand{\style@custom}{%
    \style@academic
}

\makeatother

% ======================================================================
% B. CODE BLOCKS
% ======================================================================

% codeblock environment: \begin{codeblock}[python] ... \end{codeblock}
% Light-themed code block with left accent border (uses tcolorbox + listings)
\newtcblisting{codeblock}[1][python]{
    colback=CodeBg,
    colframe=CodeBg,
    listing only,
    listing options={language=#1},
    arc=2pt,
    boxrule=0pt,
    borderline west={2.5pt}{0pt}{CodeAccent},
    left=8pt, right=8pt, top=4pt, bottom=4pt,
    before skip=8pt,
    after skip=8pt,
}

% inline code — rounded background pill
\newtcbox{\inlinecode}{
    on line,
    colback=\CodeBgColor,
    colframe=\CodeBgColor,
    boxrule=0pt,
    arc=3pt,
    boxsep=0pt,
    left=4pt, right=4pt, top=1.5pt, bottom=1.5pt,
    fontupper=\small\ttfamily,
}

% ======================================================================
% D. CALLOUT / NOTE BOXES
% ======================================================================

% Usage: \begin{bluebox}...\end{bluebox}
%        \begin{bluebox}[Custom Title]...\end{bluebox}

\definecolor{BoxBlue}{HTML}{4a90d9}
\definecolor{BoxAmber}{HTML}{e6a117}
\definecolor{BoxGreen}{HTML}{3a9a5c}
\definecolor{BoxRed}{HTML}{d94452}

\newtcolorbox{bluebox}[1][Note]{
    colback=BoxBlue!6, colframe=BoxBlue, coltitle=white,
    title={\smash{\textbf{#1}}\vphantom{N}}, fonttitle=\sffamily,
    boxrule=0.5pt, arc=2pt,
    left=6pt, right=6pt, top=4pt, bottom=4pt,
}

\newtcolorbox{amberbox}[1][Warning]{
    colback=BoxAmber!8, colframe=BoxAmber, coltitle=white,
    title={\smash{\textbf{#1}}\vphantom{N}}, fonttitle=\sffamily,
    boxrule=0.5pt, arc=2pt,
    left=6pt, right=6pt, top=4pt, bottom=4pt,
}

\newtcolorbox{greenbox}[1][Tip]{
    colback=BoxGreen!6, colframe=BoxGreen, coltitle=white,
    title={\smash{\textbf{#1}}\vphantom{N}}, fonttitle=\sffamily,
    boxrule=0.5pt, arc=2pt,
    left=6pt, right=6pt, top=4pt, bottom=4pt,
}

\newtcolorbox{redbox}[1][Danger]{
    colback=BoxRed!6, colframe=BoxRed, coltitle=white,
    title={\smash{\textbf{#1}}\vphantom{N}}, fonttitle=\sffamily,
    boxrule=0.5pt, arc=2pt,
    left=6pt, right=6pt, top=4pt, bottom=4pt,
}

% ======================================================================
% E. THEOREM ENVIRONMENTS
% ======================================================================

\newtheorem{theorem}{Theorem}[section]
\newtheorem{definition}{Definition}[section]
\newtheorem{lemma}{Lemma}[section]
\newtheorem{corollary}{Corollary}[section]
\newtheorem{proposition}{Proposition}[section]

% ======================================================================
% F. ALGORITHM (algorithm2e defaults set in preamble)
% ======================================================================

% ======================================================================
% G. TODO
% ======================================================================

% Margin todo note
\newcommand{\todox}[1]{\todo[inline, color=yellow!40]{#1}}

% Checkbox-style todo
\newcommand{\todomark}[1]{%
    $\square$~{#1}%
}
\newcommand{\tododone}[1]{%
    $\boxtimes$~{\sout{#1}}%
}

% ======================================================================
% H. MULTI-COLUMN LAYOUTS
% ======================================================================

\newenvironment{twocol}{%
    \begin{multicols}{2}%
}{%
    \end{multicols}%
}

\newenvironment{threecol}{%
    \begin{multicols}{3}%
}{%
    \end{multicols}%
}

% ======================================================================
% I. BLOCKQUOTE
% ======================================================================

\newtcolorbox{blockquote}{
    blanker,
    borderline west={3pt}{0pt}{gray!60},
    left=12pt,
    top=4pt,
    bottom=4pt,
    before skip=8pt,
    after skip=8pt,
    fontupper=\itshape\color{gray!80!black},
}

% ======================================================================
% J. SHADOW BOX (generic wrapper: border + drop shadow)
% ======================================================================

% Usage: \shadowbox{\includegraphics[...]{...}}
%        \shadowbox{any content}
% Shadow is rendered via tikz [overlay] so it does NOT affect bounding box,
% meaning centering, spacing, and caption distance stay identical to bare content.
%
% --- CONFIG (tweak these) ---
\newcommand{\ShadowColor}{gray}        % shadow color
\newcommand{\ShadowBorderColor}{gray}  % border color
\newcommand{\ShadowBorderWidth}{0.4pt} % border line width
\newcommand{\ShadowLayers}{5}          % number of blur layers (more = smoother)
\newcommand{\ShadowXShift}{0.15mm}     % base X offset per layer (positive = right)
\newcommand{\ShadowYShift}{-0.15mm}    % base Y offset per layer (negative = down)
\newcommand{\ShadowSpread}{0.12mm}     % extra spread per layer (controls blur radius)
\newcommand{\ShadowBaseOpacity}{0.12}  % opacity of innermost shadow layer
\newcommand{\ShadowFadeRate}{0.018}    % opacity decrease per layer
% ---
\newcommand{\shadowbox}[1]{%
    \begin{tikzpicture}
        \node[inner sep=0pt] (C) {#1};
        \draw[\ShadowBorderColor, line width=\ShadowBorderWidth]
            (C.south west) rectangle (C.north east);
        \begin{scope}[on background layer, overlay]
            \foreach \s in {1,...,\ShadowLayers} {
                \pgfmathsetmacro{\op}{\ShadowBaseOpacity - \ShadowFadeRate*\s}
                \fill[\ShadowColor, opacity=\op]
                    ([shift={({\ShadowXShift+\ShadowSpread*\s},{\ShadowYShift-\ShadowSpread*\s})}]C.south west)
                    rectangle
                    ([shift={({\ShadowXShift+\ShadowSpread*\s},{\ShadowYShift-\ShadowSpread*\s})}]C.north east);
            }
        \end{scope}
    \end{tikzpicture}%
}

% ======================================================================
% K. MARGIN NOTE
% ======================================================================

\newcommand{\marginmark}[1]{%
    \marginnote{\footnotesize\color{gray}#1}%
}
